\documentclass[10pt,a4paper]{amsart}
\usepackage[margin=19mm]{geometry}
\usepackage{amsmath,amssymb,mathtools}
\usepackage[scr=rsfs,cal=euler]{mathalfa}
\usepackage{xfrac}
\usepackage[unicode,hidelinks,bookmarks=false]{hyperref}
\newcommand{\ie}{i.e.\ }
\newcommand{\cf}{cf.\ }
\newcommand{\resp}{respectively\ }
\newcommand{\Subsection}{\subsection}
\providecommand{\nobreakdash}{\nobreak\hbox{-}}
\setlength{\emergencystretch}{2em}
\multlinegap0pt
\usepackage{bm}
\usepackage{bbm}
\usepackage{dsfont}
\usepackage{xspace}
\usepackage{cancel}
\usepackage{mathtools}
\usepackage{cleveref}
\usepackage{amsthm}
\makeatletter
\@ifundefined{theorem}{\newtheorem{theorem}{Theorem}}{}
\@ifundefined{lemma}{\newtheorem{lemma}[theorem]{Lemma}}{}
\makeatother
\DeclareMathOperator{\Aut}{Aut}
\DeclareMathOperator{\St}{St}
\DeclareMathOperator{\pr}{pr}
\newcommand{\N}{\mathbb{N}}
\newcommand{\T}{\mathcal{T}}
\newcommand{\absH}[1]{\Vert #1 \Vert_H}
\newcommand{\bigAbsH}[1]{\bigl\Vert #1 \bigr\Vert_H}
\renewcommand\leq\leqslant
\renewcommand\geq\geqslant
\title[A branch group with unsolvable conjugacy problem]{A branch group with unsolvable conjugacy problem}
\author{Alex Bishop, Eduard Schesler}
\date{Source: arXiv:2509.12161v2 (CC BY 4.0)}
\begin{document}
\maketitle

\noindent\emph{Source and licence.} A branch group with unsolvable conjugacy problem. Version arXiv:2509.12161v2 (CC BY 4.0), distributed under the Creative Commons Attribution 4.0 International licence, as recorded in the arXiv licence metadata for this exact version and shown on the abstract page https://arxiv.org/abs/2509.12161v2. Full rights notice: licenses/arxiv-2509.12161-CC-BY-4.0.txt.\par\smallskip

\noindent\textit{Editorial scope.} Counted unit: the Theorem whose source label is \texttt{thm:fratini-embedding} in the article (source0.tex lines 759--763, source label \texttt{thm:fratini-embedding}), delivered together with the article's own complete original proof of it (source0.tex lines 765--931, one \texttt{proof} environment of 167 lines). No mathematics is written, repaired or re-derived: statement and proof are the article's own text line for line. Required context retained uncounted: eq:section-formula (lines 356--358); lem:length-contraction-base-case (lines 674--694); lem:length-contraction (lines 711--723). Every definition, display and label of those lines is retained unchanged. Omitted from this package and not redistributed: 1--355, 359--673, 695--710, 724--758, 764--764, 932--1110. Nothing inside a retained excerpt is omitted or altered: every retained body is delivered exactly as the article's lines are written, so no omission receipt of any kind accompanies this package. Citations: the retained excerpts carry no citation command at all, so no bibliography entry of the article is transcribed and no reference is redistributed. Reference adaptation: none was needed and none was made; every reference inside the delivered unit resolves inside it, so no unresolved reference or citation is produced. Printed slips kept literal: no printed word, symbol, hypothesis or number of the counted statement or of its proof was altered. Layout and class: the article's own class is \texttt{amsart} with packages including inputenc, mathtools, amssymb, amsfonts, amsmath, enumerate, graphicx, verbatim, tikz-cd, comment, tikz, enumitem, yhmath, thmtools; the wrapper is the lane's pinned \texttt{amsart} editorial wrapper at 10pt on A4, which restates the theorem-like environments the retained text needs and adds the article's own macro definitions that the retained text needs (\texttt{\textbackslash Aut}, \texttt{\textbackslash N}, \texttt{\textbackslash St}, \texttt{\textbackslash T}, \texttt{\textbackslash absH}, \texttt{\textbackslash bigAbsH}, \texttt{\textbackslash leq}, \texttt{\textbackslash pr}), with extra wrapper packages (bm, bbm, dsfont, xspace, cancel, mathtools, cleveref). The delivered numbering is the unit's own. Assets: no retained span contains a figure environment, a graphics-inclusion command, a PDF-page inclusion, a table or a TikZ picture; no image file is copied into this package, the package declares no asset dependency and no image file is redistributed with it. Licence: version arXiv:2509.12161v2 of this article is distributed under the Creative Commons Attribution 4.0 International licence, as recorded in the arXiv licence metadata for this exact version and shown on the abstract page https://arxiv.org/abs/2509.12161v2. A full rights notice is delivered at licenses/arxiv-2509.12161-CC-BY-4.0.txt. Characters: every retained excerpt is pure ASCII, so the delivered engine, XeLaTeX with its default Latin Modern text font, reports no missing character.
\begin{equation}\label{eq:section-formula}
(\alpha \beta)_v = \alpha_{\beta(v)} \beta_v,
\end{equation}

\begin{lemma}\label{lem:length-contraction-base-case}
  Let $\ell \in \N_0$, and suppose that $w, w' \in F_\ell$ have the forms
  \[
    w = b_1 h_1 b_2
    \qquad\text{and}\qquad
    w' = b_1 h_1 b_2 h_2 b_3.
  \]
  We write $\gamma = \pr_\ell(w)$ and $\gamma'=\pr_\ell(w')$ for the element of $\Gamma_\ell$ corresponding to $w$ and $w'$, respectively.
  Then, for each $d\in X_\ell$, there are elements $w_d,w_d' \in F_{\ell+1}$ with
  \begin{align*}
    w_d &\in
        B_{\ell+1}
        \cup \{ h_1 \}\text{ and}\\
    w'_d &\in
        B_{\ell+1}
        \cup \{ h_1 h_2 \}
        \cup \{ h_1 \beta \mid \beta \in B_{\ell+1} \}
        \cup \{ \beta h_2 \mid \beta \in B_{\ell+1} \}
  \end{align*}
  such that $\gamma_d = \pr_{\ell + 1}(w_d)$ and $\gamma'_d = \pr_{\ell+1}(w'_d)$.
\end{lemma}

\begin{lemma}\label{lem:length-contraction}
Let $\ell \in \N_0$, let $w \in F_{\ell}$ be a word of the form
\[
w = b_1 h_1 b_2 h_2 \cdots b_n h_n b_{n+1},
\]
where $b_i \in B_\ell$ and $h_i\in H$, and let $\gamma = \pr_\ell(w)$.
Then for each $d \in X_{\ell+1}$, we may find a word of the form
\[
w_d = \beta_1 h'_1 \beta_2 h'_2 \cdots \beta_m h'_m \beta_{m+1} \in F_{\ell+1}
\]
such that $\gamma_d = \pr_{\ell+1}(w_d)$ with $m \leq \lceil n/2 \rceil$, $\beta_i \in B_{\ell+1}$, and $h_i'\in H^*$ such that $h_1' h_2' \cdots h_m'\in H^*$ is a fragmented subword of $h_1 h_2 \cdots h_n \in H^*$.
In particular we have $\bigAbsH{\gamma_d} \leq \lceil \absH{\gamma}/2 \rceil$  for each $d\in X_{\ell+1}$.
\end{lemma}

\begin{theorem}\label{thm:fratini-embedding}
Let $\ell \in \N_0$.
Two elements $g,k \in G$ are conjugate in $G$ if and only if the corresponding elements $\widetilde{g}^{[\ell]},\widetilde{k}^{[\ell]}\in \Gamma_\ell$ are conjugate in $\Gamma_\ell$.
  That is, for every $\ell\in \N_0$, the map $g\mapsto \widetilde{g}^{[\ell]}$ is a Frattini embedding of $G$ into $\Gamma_\ell$.
\end{theorem}

\begin{proof}
Since the map $g\mapsto \widetilde{g}^{[\ell]}$ is a homomorphism it is clear that $\widetilde{g}^{[\ell]},\widetilde{k}^{[\ell]}\in \Gamma_\ell$ are conjugate in $\Gamma_\ell$ if $g,k \in G$ are conjugate in $G$.
Suppose now that $g,k\in G$ are elements that are not conjugate in $G$.
Thus, to complete our proof, it remains to be shown that $\widetilde{g}^{[\ell]}$ and $\widetilde{k}^{[\ell]}$ are not conjugate in $\Gamma_\ell$.
Since $g\mapsto \widetilde{g}^{[\ell]}$ is injective, this property immediately follows if one of $g$ or $k$ is the identity.
In the remainder of this proof, we assume without loss of generality that neither $g$ nor $k$ is the identity of $G$.			 
Suppose for contradiction, that $\widetilde{g}^{[\ell]}$ and $\widetilde{k}^{[\ell]}$ are conjugate, that is, that there exists some $\gamma\in \Gamma_\ell$ for which $\gamma^{-1} \widetilde{g}^{[\ell]} \gamma = \widetilde{k}^{[\ell]}$.
We will show, by induction on $\absH{\gamma}$, that no such element $\gamma\in \Gamma_\ell$ can exist.
We begin with the base cases $\absH{\gamma} = 0$ and $\absH{\gamma}=1$, as follows.

\medskip

\noindent\underline{Base case}: $\absH{\gamma} = 0$.
\nopagebreak

\nopagebreak\smallskip
\nopagebreak\noindent
To begin, suppose that $\absH{\gamma}=0$, that is, $\gamma\in B_\ell$.

Recall from the definition of $\widetilde{\cdot}^{[\ell]}$ that
$(\widetilde{g}^{[\ell]})_{x_{\ell+1}} = \widetilde{g}^{[\ell+1]}$ and
$(\widetilde{k}^{[\ell]})_{x_{\ell+1}} = \widetilde{k}^{[\ell+1]}$, neither of which are finitary;
and that for each $a\in X_{\ell+1}\setminus\{x_{\ell+1}\}$, we have $(\widetilde{g}^{[\ell]})_{a} \in B_{\ell+1}$ and
$(\widetilde{k}^{[\ell]})_{a} \in B_{\ell+1}$, which are finitary.
Moreover, notice that, if $\gamma^{-1} \widetilde{g}^{[\ell]} \gamma = \widetilde{k}^{[\ell]}$, then
\[
(\gamma^{-1} \widetilde{g}^{[\ell]} \gamma)_{x_{\ell+1}}
= \widetilde{k}^{[\ell+1]}.
\]
Hence, it must be the case that $\gamma\cdot x_{\ell+1}=x_{\ell+1}$ and thus $\widetilde{g}^{[\ell+1]} = \widetilde{k}^{[\ell+1]}$.
Since the map $\widetilde{\cdot}^{[\ell+1]} \colon G \to \Aut(\T_X^{[\ell+1]})$ is injective, this would imply that $g=k$ which contradicts our assumption that $g$ and $k$ are not conjugate in $G$.

\medskip

\noindent\underline{Base case}: $\absH{\gamma} = 1$.
\nopagebreak

\nopagebreak\smallskip
\nopagebreak\noindent
Suppose that $\absH{\gamma}=1$, then there exists an element $w\in F_\ell$ with $\gamma = \pr_\ell(w)$ and
\begin{equation}\label{eq:base-case1}
    w = b_1 h_1 b_2\in F_\ell,
\end{equation}
where $b_i\in B_\ell$ and $h_1 \in H$.
Since $\widetilde{g}^{[\ell]}\in \St_{\Gamma_\ell}(1)$ and $\gamma^{-1} \widetilde{g}^{[\ell]} \gamma = \widetilde{k}^{[\ell]}$, there exist $a,b\in X_{\ell+1}$ such that
\begin{equation}\label{eq:eq1}
    \bigl(\gamma_a\bigr)^{-1}
    \cdot (\widetilde{g}^{[\ell]})_b
    \cdot \gamma_a
    =
    (\widetilde{k}^{[\ell]})_{x_{\ell+1}} = \widetilde{k}^{[\ell+1]}.
\end{equation}
From the definition of $\widetilde{g}^{[\ell]}$, we then see that either $(\widetilde{g}^{[\ell]})_b \in B_{\ell+1}$ or $(\widetilde{g}^{[\ell]})_b = \widetilde{g}^{[\ell+1]}$.
Further, from \cref{lem:length-contraction-base-case}, we see that \[\gamma_{a} \in B_{\ell+1} \cup \left\{\widetilde{h_1}^{[\ell+1]}\right\}.\]
From (\ref{eq:eq1}), the remainder of the proof of this base case falls into 4 cases as follows.

\medskip

\noindent\textit{Case 1:}
$\tau^{-1} \cdot \widetilde{g}^{[\ell+1]}\cdot \tau = \widetilde{k}^{[\ell+1]}$ for some $\tau\in B_{\ell+1}$.

\smallskip\noindent
In this case, our contradiction follows from the previous base case.

\medskip

\noindent\textit{Case 2:}
$\tau^{-1}\cdot \sigma\cdot \tau = \widetilde{k}^{[\ell+1]}$ for some $\sigma,\tau\in B_{\ell+1}$.

\smallskip\noindent
In this case, the contradiction follows as $\widetilde{k}^{[\ell+1]}$ is not finitary and in particular does not belong to $B_{\ell+1}$ for $k\neq 1$.

\medskip

\noindent\textit{Case 3:}
$(\widetilde{h_1}^{[\ell+1]})^{-1}\cdot (\widetilde{g}^{[\ell+1]})\cdot (\widetilde{h_1}^{[\ell+1]}) = \widetilde{k}^{[\ell+1]}$.

\smallskip\noindent
Since $\widetilde{\cdot}^{[\ell+1]} \colon H\to\Gamma_{\ell+1}$ is injective, it would then follow that $h_1^{-1} g h_1 = k$ in $H=G\times A$.
Since $H$ is a direct product of $G$ and $A$, with $h_1 = (x,y)\in G\times A = H$, we see that this implies that $x^{-1} g x = k$ in $G$.
This contradicts our assumption that $g$ and $k$ are not conjugate in $G$.

\medskip

\noindent\textit{Case 4:}
$(\widetilde{h_1}^{[\ell+1]})^{-1}\cdot \sigma\cdot (\widetilde{h_1}^{[\ell+1]}) = \widetilde{k}^{[\ell+1]}$ for some $\sigma\in B_{\ell+1}$.

\smallskip\noindent
If $\sigma = 1$, then $\widetilde{k}^{[\ell+1]} = 1$, which would contradict our assumption that $k$ is non-trivial.
Suppose therefore that $\sigma \in B_{\ell+1}\setminus \{1\}$.
Then, a contradiction arises as
\[									 
(\widetilde{h_1}^{[\ell+1]})^{-1}\cdot \sigma\cdot (\widetilde{h_1}^{[\ell+1]})\notin \St_{\Gamma_{\ell+1}}(1)
\qquad
\text{while}
\qquad
\widetilde{k}^{[\ell+1]}\in \St_{\Gamma_{\ell+1}}(1).
\]
This completes all subcases of the base case $\absH{\gamma}=1$.

\medskip

\noindent\underline{Inductive step}:
\nopagebreak

\nopagebreak\smallskip
\nopagebreak\noindent
We now assume for induction that there is some $m\in \N$ such that for every $\ell\in\N_0$ there is no word $\gamma'\in\Gamma_\ell$ with $\absH{\gamma'}\leq m$ such that $(\gamma')^{-1} \widetilde{g}^{[\ell]} \gamma' = \widetilde{k}^{[\ell]}$.
Suppose for contradiction that there exists $\gamma\in \Gamma_\ell$ with $\absH{\gamma}=m+1$ for which $\gamma^{-1} \widetilde{g}^{[\ell]} \gamma = \widetilde{k}^{[\ell]}$.
By the definition of $\absH{\cdot}$, there must exist some element $w\in F_\ell$ which can be written as
\[
    w = b_1 h_1 b_2 h_2 \cdots b_m h_m b_{m+1} h_{m+1} b_{m+2} 
\]
where $b_i\in B_\ell$, $h_i\in H$, and $\gamma = \pr_\ell(w)$.

Notice that if $\gamma^{-1} \widetilde{g}^{[\ell]} \gamma = \widetilde{k}^{[\ell]}$, then it would follow that

\begin{equation}\label{eq:projections}
    \bigl(\gamma^{-1} \widetilde{g}^{[\ell]} \gamma\bigr)_{x_{\ell+1}}
    =
    (\widetilde{k}^{[\ell]})_{x_{\ell+1}}
    =
    \widetilde{k}^{[\ell+1]}.
\end{equation}
Let $\beta\in B_\ell$ be the unique element for which
$
    \gamma\beta\in \St_{\Gamma_\ell}(1).
$
From \eqref{eq:section-formula}, \eqref{eq:projections} and the definition of $\widetilde{\cdot}^{[\ell]}\colon H \to \Gamma_\ell$, we see that there exists some $d\in X_{\ell+1}$ for which
\begin{equation}\label{eq:projections2}
    \bigl((\gamma\beta)_{d}\bigr)^{-1}
    \cdot (\widetilde{g}^{[\ell]})_d
    \cdot (\gamma\beta)_{d}
    =
    \widetilde{k}^{[\ell+1]}.
\end{equation}
Note that we either have $(\widetilde{g}^{[\ell]})_d\in B_\ell$ or $(\widetilde{g}^{[\ell]})_d = \widetilde{g}^{[\ell+1]}$.
We consider these two possibilities in separate cases.

\medskip

\noindent\textit{Case 1}: $(\widetilde{g}^{[\ell]})_d\in B_\ell$.
\nopagebreak

\nopagebreak\smallskip
\nopagebreak\noindent
Same arguments as in case 4 of the proof of the base case $\absH{\gamma}=1$.

\medskip

\noindent\textit{Case 2}: $(\widetilde{g}^{[\ell]})_d = \widetilde{g}^{[\ell+1]}$.
\nopagebreak

\nopagebreak\smallskip
\nopagebreak\noindent
From \cref{lem:length-contraction}, we see that $\bigAbsH{(\gamma\beta)_d}< \absH{\gamma}$.
Thus, we have a $\gamma'\in \Gamma_{\ell+1}$ with $\absH{\gamma'}< \absH{\gamma}$ for which $\gamma'^{-1} \widetilde{g}^{[\ell+1]}\gamma' = \widetilde{k}^{[\ell+1]}$ which contradicts our induction hypothesis.

\medskip

\noindent\underline{Conclusion}:
\nopagebreak

\nopagebreak\smallskip
\nopagebreak\noindent
In all cases, we see that a contradiction arises if we assume that there exists some $\gamma\in \Gamma_{\ell}$ for which $\gamma^{-1} \widetilde{g}^{[\ell]}\gamma = \widetilde{k}^{[\ell]}$ where $g$ and $k$ are not conjugate in $G$.
\end{proof}

\end{document}
